% GENERATED FILE. Edit canonical lessons, then run npm run generate:books.
% canonical-source-hash: fnv1a64:607caa082c533c4f
% canonical-lessons: JA-C63-yume, JA-C63-matsuri, JA-C63-nishi, JA-C63-higashi, JA-C63-minami

\chapter{Dream, Festival, West, East, South}
\label{ch:ja-dream-festival-west-east-south}
\hlchaptermodality{63}

\begin{chapteropening}
\textbf{By the end of this chapter:} \emph{I can say a dream, a festival, west, east and south.}

Complete the last lesson of chapter 63: i can say a dream, a festival, west, east and south.
\end{chapteropening}

\section[yume]{\ja{ゆめ} (yume) — a dream}

\label{lesson:JA-C63-yume}



\begin{quote}
Before the new one: say the Japanese for a colour, then the Japanese for a shape.
\end{quote}

\subsection*{You'll want to know: \ja{ゆめ}}
\textbf{\ja{ゆめ}} — \emph{yume} — \textquotedblleft{}a dream\textquotedblright{}.

\begin{grammarlens}[title={What you've built}]
One more word for this chapter.
\end{grammarlens}

\subsection*{Your turn}
\begin{itemize}
\raggedright
  \item \emph{Say it:} \emph{yume}
  \item \emph{Say it:} \emph{yume}, once more
  \item \emph{Say it:} say \emph{yume}
  \item \emph{Recall:} say the Japanese for a colour, then the Japanese for a shape, then say \emph{yume} again
\end{itemize}

\begin{tcolorbox}[breakable,colback=teal!4,colframe=teal!35!black,title={Before you move on}]
What does \ja{ゆめ} mean? (\textquotedblleft{}A dream\textquotedblright{}.) Say it once more.
\end{tcolorbox}
\section[matsuri]{\ja{まつり} (matsuri) — a festival}

\label{lesson:JA-C63-matsuri}



\begin{quote}
Before the new one: say the Japanese for a shape, then the Japanese for a dream.
\end{quote}

\subsection*{You'll want to know: \ja{まつり}}
\textbf{\ja{まつり}} — \emph{matsuri} — \textquotedblleft{}a festival\textquotedblright{}.

\begin{grammarlens}[title={What you've built}]
One more word for this chapter.
\end{grammarlens}

\subsection*{Your turn}
\begin{itemize}
\raggedright
  \item \emph{Say it:} \emph{matsuri}
  \item \emph{Say it:} \emph{matsuri}, once more
  \item \emph{Say it:} say \emph{matsuri}
  \item \emph{Recall:} say the Japanese for a shape, then the Japanese for a dream, then say \emph{matsuri} again
\end{itemize}

\begin{tcolorbox}[breakable,colback=teal!4,colframe=teal!35!black,title={Before you move on}]
What does \ja{まつり} mean? (\textquotedblleft{}A festival\textquotedblright{}.) Say it once more.
\end{tcolorbox}
\section[nishi]{\ja{にし} (nishi) — west}

\label{lesson:JA-C63-nishi}



\begin{quote}
Before the new one: say the Japanese for a dream, then the Japanese for a festival.
\end{quote}

\subsection*{You'll want to know: \ja{にし}}
\textbf{\ja{にし}} — \emph{nishi} — \textquotedblleft{}west\textquotedblright{}.

\begin{grammarlens}[title={What you've built}]
One more word for this chapter.
\end{grammarlens}

\subsection*{Your turn}
\begin{itemize}
\raggedright
  \item \emph{Say it:} \emph{nishi}
  \item \emph{Say it:} \emph{nishi}, once more
  \item \emph{Say it:} say \emph{nishi}
  \item \emph{Recall:} say the Japanese for a dream, then the Japanese for a festival, then say \emph{nishi} again
\end{itemize}

\begin{tcolorbox}[breakable,colback=teal!4,colframe=teal!35!black,title={Before you move on}]
What does \ja{にし} mean? (\textquotedblleft{}West\textquotedblright{}.) Say it once more.
\end{tcolorbox}
\section[higashi]{\ja{ひがし} (higashi) — east}

\label{lesson:JA-C63-higashi}



\begin{quote}
Before the new one: say the Japanese for a festival, then the Japanese for west.
\end{quote}

\subsection*{You'll want to know: \ja{ひがし}}
\textbf{\ja{ひがし}} — \emph{higashi} — \textquotedblleft{}east\textquotedblright{}.

\begin{grammarlens}[title={What you've built}]
One more word for this chapter.
\end{grammarlens}

\subsection*{Your turn}
\begin{itemize}
\raggedright
  \item \emph{Say it:} \emph{higashi}
  \item \emph{Say it:} \emph{higashi}, once more
  \item \emph{Say it:} say \emph{higashi}
  \item \emph{Recall:} say the Japanese for a festival, then the Japanese for west, then say \emph{higashi} again
\end{itemize}

\begin{tcolorbox}[breakable,colback=teal!4,colframe=teal!35!black,title={Before you move on}]
What does \ja{ひがし} mean? (\textquotedblleft{}East\textquotedblright{}.) Say it once more.
\end{tcolorbox}
\section[minami]{\ja{みなみ} (minami) — south}

\label{lesson:JA-C63-minami}



\begin{quote}
Before the new one: say the Japanese for west, then the Japanese for east.
\end{quote}

\subsection*{You'll want to know: \ja{みなみ}}
\textbf{\ja{みなみ}} — \emph{minami} — \textquotedblleft{}south\textquotedblright{}.

\begin{grammarlens}[title={What you've built}]
One more word for this chapter.
\end{grammarlens}

\subsection*{Your turn}
\begin{itemize}
\raggedright
  \item \emph{Say it:} \emph{minami}
  \item \emph{Say it:} \emph{minami}, once more
  \item \emph{Say it:} say \emph{minami}
  \item \emph{Recall:} say the Japanese for west, then the Japanese for east, then say \emph{minami} again
\end{itemize}

\begin{tcolorbox}[breakable,colback=teal!4,colframe=teal!35!black,title={Before you move on}]
What does \ja{みなみ} mean? (\textquotedblleft{}South\textquotedblright{}.) Say it once more.
\end{tcolorbox}
